\documentclass[11pt,a4paper]{article}

% ---------------------------------------------------------------------------
%  Ciphers -- Security Report
% ---------------------------------------------------------------------------
\usepackage[utf8]{inputenc}
\usepackage[T1]{fontenc}
\usepackage{lmodern}
\usepackage{amsmath}
\usepackage{amssymb}
\usepackage[margin=1in]{geometry}
\usepackage{microtype}
\usepackage{booktabs}
\usepackage{longtable}
\usepackage{enumitem}
\usepackage{xcolor}
\usepackage{titlesec}
\usepackage{fancyhdr}
\usepackage{hyperref}

\definecolor{ink}{HTML}{0E1B2B}
\definecolor{accent}{HTML}{0E7C86}
\definecolor{good}{HTML}{1B7A2E}
\definecolor{warn}{HTML}{B23A48}

\hypersetup{
  colorlinks=true,
  linkcolor=accent,
  urlcolor=accent,
  citecolor=accent,
  pdftitle={Ciphers Security Report},
  pdfauthor={Jean-Francois Lachance-Caumartin}
}

\titleformat{\section}{\Large\bfseries\color{ink}}{\thesection}{0.6em}{}
\titleformat{\subsection}{\large\bfseries\color{ink}}{\thesubsection}{0.6em}{}

\pagestyle{fancy}
\fancyhf{}
\lhead{\textsc{Ciphers} — Security Report}
\rhead{\thepage}
\renewcommand{\headrulewidth}{0.4pt}

\newcommand{\yes}{\textcolor{good}{\textbf{Yes}}}
\newcommand{\no}{\textcolor{warn}{\textbf{No}}}
\newcommand{\code}[1]{\texttt{#1}}

% ---------------------------------------------------------------------------
\begin{document}

\begin{titlepage}
\centering
\vspace*{3cm}
{\Huge\bfseries\color{ink} Ciphers}\\[0.4cm]
{\Large Security Report}\\[1.5cm]
{\large Authenticated file encryption with a post-quantum hybrid KEM}\\[0.3cm]
{\large Version 1.0.4}\\[2.5cm]

\begin{tabular}{rl}
\textbf{Author:} & Jean-François Lachance-Caumartin\\[0.2cm]
\textbf{Document:} & Security architecture \& threat model\\[0.2cm]
\textbf{Status:} & Public\\
\end{tabular}

\vfill
{\large \today}
\end{titlepage}

\tableofcontents
\newpage

% ---------------------------------------------------------------------------
\section{Overview}

Ciphers is a small GTK3 desktop application for Linux that encrypts and
decrypts arbitrary files under a single user password. It provides modern
authenticated encryption (AEAD) and an optional \emph{post-quantum hybrid key
encapsulation} layer that combines CRYSTALS-Kyber-1024 with the X448
elliptic curve.

This document describes the cryptographic design, the on-disk file format,
the memory-hardening measures, the threat model, and — equally important —
the properties that Ciphers explicitly does \emph{not} provide. It is
intended for users evaluating the tool for sensitive data and for reviewers
auditing its guarantees.

\subsection{Design goals}
\begin{itemize}[leftmargin=1.4em]
  \item Confidentiality and integrity of file contents under a password.
  \item Optional protection against a future quantum adversary
        (``harvest-now, decrypt-later'').
  \item Resistance to tampering, reordering and truncation of ciphertext.
  \item Keeping secrets (password, derived keys, plaintext) off disk.
  \item A small, auditable codebase built on vetted libraries.
\end{itemize}

\subsection{Non-goals}
\begin{itemize}[leftmargin=1.4em]
  \item Hiding the existence, size or type of an encrypted file
        (no steganography or plausible deniability).
  \item Filename, metadata or directory-structure encryption.
  \item Multi-recipient / public-key sharing workflows.
  \item Protection against a compromised host (keylogger, malicious GTK
        input method, cold-boot attack on live RAM).
\end{itemize}

% ---------------------------------------------------------------------------
\section{Cryptographic primitives}

All symmetric primitives are provided by \textbf{libsodium}; the X448 curve is
provided by \textbf{OpenSSL} (\code{libcrypto}); the Kyber-1024 implementation
is the CRYSTALS reference code (\code{KYBER\_K=4}, SHAKE/FIPS-202 variant),
with randomness drawn from libsodium. Argon2id is provided by
\textbf{libargon2}.

\begin{center}
\begin{tabular}{@{}lll@{}}
\toprule
\textbf{Purpose} & \textbf{Primitive} & \textbf{Parameters}\\
\midrule
AEAD (default)      & AES-256-GCM              & 256-bit key, 96-bit nonce, 128-bit tag\\
AEAD (alt.)         & XChaCha20-Poly1305       & 256-bit key, 192-bit nonce, 128-bit tag\\
AEAD (alt.)         & ChaCha20-Poly1305 (IETF) & 256-bit key, 96-bit nonce, 128-bit tag\\
Key derivation      & Argon2id                 & see Table~\ref{tab:kdf}\\
PQ KEM              & CRYSTALS-Kyber-1024      & NIST security level 5\\
Classical KEM       & X448 ECDH                & 448-bit curve\\
Hybrid combiner     & BLAKE2b (keyed)          & domain ``\code{HYBRID}'', 32-byte output\\
Key wrap            & XChaCha20-Poly1305       & wraps the hybrid secret key\\
\bottomrule
\end{tabular}
\end{center}

\subsection{Key derivation}
\label{sec:kdf}

The password is stretched with Argon2id. Three presets are offered; the
parameters are stored in the file header so decryption reproduces them.

\begin{table}[h]
\centering
\begin{tabular}{@{}llll@{}}
\toprule
\textbf{Preset} & \textbf{Memory} & \textbf{Iterations} & \textbf{Parallelism}\\
\midrule
Basic   & 256\,MiB & 3 & 4 lanes\\
Medium\,$^\dagger$ & 1\,GiB   & 3 & 4 lanes\\
Strong  & 4\,GiB   & 4 & 8 lanes\\
\bottomrule
\end{tabular}
\caption{Argon2id presets. $^\dagger$Medium is the minimum recommended for
sensitive data.}
\label{tab:kdf}
\end{table}

A fresh 16-byte random salt is generated per file, so identical passwords on
different files yield independent keys and the KDF output cannot be
precomputed or shared across files.

% ---------------------------------------------------------------------------
\section{File format}

Every encrypted file begins with a fixed 40-byte header:

\begin{center}
\begin{tabular}{@{}rll@{}}
\toprule
\textbf{Offset} & \textbf{Size} & \textbf{Field}\\
\midrule
0  & 8  & Magic \code{"CIPHERS\textbackslash0"}\\
8  & 1  & Format version (1 = password-only, 2 = hybrid)\\
9  & 1  & Cipher id\\
10 & 1  & KDF id (1 = Argon2id)\\
11 & 1  & KDF level (informational)\\
12 & 4  & Argon2 iterations (little-endian)\\
16 & 4  & Argon2 memory in KiB (little-endian)\\
20 & 4  & Argon2 parallelism (little-endian)\\
24 & 16 & Salt\\
\bottomrule
\end{tabular}
\end{center}

In hybrid files (version 2) a \emph{hybrid block} follows the header,
containing the wrap nonce, the wrapped hybrid secret key and the KEM
ciphertext. The base AEAD nonce follows, then the payload frames.

\subsection{Chunked authenticated streaming}

The plaintext is split into 64\,KiB chunks. Each chunk is encrypted as an
independent AEAD frame:
\[
  \text{frame} = \underbrace{\texttt{len}}_{4\,\text{B}} \;\Vert\;
  \underbrace{\text{ciphertext} \,\Vert\, \text{tag}}_{\texttt{len}\,\text{B}}.
\]
The per-chunk nonce is the random base nonce with a 64-bit little-endian
counter XORed into its trailing eight bytes, guaranteeing a unique nonce per
chunk under a given key. The associated data bound into each frame is
\[
  \text{AD} = \underbrace{\texttt{counter}}_{8\,\text{B}} \;\Vert\;
  \underbrace{\texttt{final\_flag}}_{1\,\text{B}}.
\]
Binding the counter authenticates chunk \emph{ordering}; binding the
end-of-stream flag authenticates that the last chunk really is the last one.
Together these detect:
\begin{itemize}[leftmargin=1.4em]
  \item \textbf{Tampering} — any modified ciphertext byte fails its Poly1305
        / GCM tag.
  \item \textbf{Reordering} — swapping frames breaks the counter binding.
  \item \textbf{Truncation} — dropping the final frame(s) leaves a chunk whose
        \code{final\_flag} no longer matches; the missing-final-block check
        also fires.
  \item \textbf{Splicing / extension} — appending frames after the
        authenticated final frame is rejected.
\end{itemize}
These properties were validated empirically with round-trip, bit-flip,
truncation and frame-swap tests across all three ciphers and the hybrid mode.

% ---------------------------------------------------------------------------
\section{Post-quantum hybrid KEM}

When hybrid mode is enabled the file's AEAD key is not derived directly from
the password. Instead:

\begin{enumerate}[leftmargin=1.6em]
  \item A fresh Kyber-1024 + X448 keypair is generated per file.
  \item The hybrid \emph{secret} key (\code{kyber\_sk} $\Vert$
        \code{x448\_sk}) is wrapped with XChaCha20-Poly1305 under a master
        key derived from the password via Argon2id, bound to the domain
        string \code{"CIPHERS-HYBRID-WRAP"}.
  \item The KEM encapsulates to the matching public key, producing a 32-byte
        shared secret that becomes the AEAD file key.
  \item The shared secret is the keyed BLAKE2b combination of the Kyber and
        X448 secrets, so the key stays secure as long as \emph{either}
        primitive holds.
\end{enumerate}

On decryption the password unwraps the secret key, which decapsulates the KEM
ciphertext back to the AEAD key. A single password therefore unlocks the
file while the construction remains secure against a quantum adversary who
can break X448 but not Kyber (and vice versa).

\paragraph{Why public keys are not stored.} Decapsulation needs only the
secret key and the KEM ciphertext, so the per-file public keys are
deliberately omitted. This removes $\sim$1.6\,KiB of otherwise unauthenticated
header bytes: the wrap nonce and wrapped secret key are covered by the wrap's
Poly1305 tag, and the KEM ciphertext is covered transitively — altering it
changes the decapsulated key and fails every subsequent frame tag.

% ---------------------------------------------------------------------------
\section{Memory hardening}

\begin{itemize}[leftmargin=1.4em]
  \item \textbf{No swap.} Passwords, derived keys and plaintext buffers are
        pinned with \code{sodium\_mlock}, which prevents them from being
        paged to the swap file and marks them \code{MADV\_DONTDUMP}.
  \item \textbf{Zeroed after use.} The corresponding \code{sodium\_munlock}
        wipes each buffer before releasing it.
  \item \textbf{No core dumps.} \code{RLIMIT\_CORE} is set to zero and, on
        Linux, \code{PR\_SET\_DUMPABLE} is cleared (blocking \code{ptrace}
        and \code{/proc}-based core capture).
  \item \textbf{Guarded password field.} The password \code{GtkEntry} is
        backed by a custom \code{GtkEntryBuffer} whose text lives in
        \code{sodium\_malloc}'d guarded memory (guard pages, locked, zeroed
        on free).
\end{itemize}

\paragraph{Residual exposure.} GTK itself may keep transient copies of typed
text in ordinary heap memory for on-screen rendering (Pango), the clipboard
and the input method. These copies are outside Ciphers' control. The
hardening therefore \emph{reduces} but cannot fully eliminate in-memory
exposure of the password.

% ---------------------------------------------------------------------------
\section{Threat model}

\subsection{In scope — what Ciphers defends against}

\begin{itemize}[leftmargin=1.4em]
  \item A passive adversary who obtains the ciphertext file and attempts to
        recover the plaintext without the password.
  \item An active adversary who modifies, reorders, truncates or extends the
        ciphertext: all such tampering is detected and decryption aborts.
  \item Offline password guessing, slowed by memory-hard Argon2id and a
        per-file salt.
  \item A future quantum adversary (hybrid mode), under the
        ``harvest-now, decrypt-later'' assumption.
  \item Recovery of secrets from the swap file or a core dump on a
        non-compromised host.
\end{itemize}

\subsection{Out of scope — what Ciphers does \emph{not} protect against}

\begin{itemize}[leftmargin=1.4em]
  \item \textbf{A compromised host.} Keyloggers, malicious input methods,
        memory scrapers running as the same or a more privileged user, or
        cold-boot attacks on live RAM are out of scope.
  \item \textbf{Weak passwords.} Argon2id slows guessing but cannot rescue a
        low-entropy password. Strength is the user's responsibility.
  \item \textbf{Metadata leakage.} The file's existence, exact size
        (to within one 64\,KiB chunk plus per-frame overhead), and the chosen
        cipher/KDF parameters are visible in the clear. Filenames are not
        encrypted.
  \item \textbf{Deniability.} The \code{"CIPHERS\textbackslash0"} magic makes
        an encrypted file trivially identifiable as such.
  \item \textbf{Long-term key escrow / recovery.} There is no backdoor and no
        recovery mechanism; a forgotten password means permanent data loss.
\end{itemize}

% ---------------------------------------------------------------------------
\section{Security properties at a glance}

\begin{center}
\begin{longtable}{@{}p{0.62\textwidth}c@{}}
\toprule
\textbf{Property} & \textbf{Provided}\\
\midrule
\endhead
Confidentiality of file contents              & \yes\\
Integrity / tamper detection (AEAD)           & \yes\\
Reordering detection                          & \yes\\
Truncation \& extension detection             & \yes\\
Per-file unique salt and nonce                & \yes\\
Memory-hard password stretching (Argon2id)    & \yes\\
Post-quantum confidentiality (hybrid mode)    & \yes\\
Secrets kept off swap / out of core dumps     & \yes\\
Authenticated KDF \& format parameters        & \yes (transitively)\\
\midrule
Filename / metadata encryption                & \no\\
Hiding that a file is encrypted (deniability) & \no\\
Protection on a compromised host              & \no\\
Password recovery / escrow                    & \no\\
Multi-recipient public-key sharing            & \no\\
\bottomrule
\end{longtable}
\end{center}

% ---------------------------------------------------------------------------
\section{Hardening notes \& known limitations}

\begin{itemize}[leftmargin=1.4em]
  \item \textbf{Untrusted headers are validated.} On decryption the Argon2id
        parameters read from the file are range-checked (iterations,
        parallelism, and the rule $m\_cost \ge 8 \times parallelism$, capped
        at 4\,GiB / 16 iterations / 16 lanes) so a malicious file cannot force
        an out-of-memory allocation or a hang.
  \item \textbf{Atomic output.} Output is written to a sibling temporary file
        and \code{rename()}d into place only on success, so a wrong password
        or a cancelled run never truncates a pre-existing output. The
        same-file check covers the input, the output and the temporary file.
  \item \textbf{AES-256-GCM availability.} AES-256-GCM requires hardware AES
        support; on CPUs without it the cipher is unavailable and the
        ChaCha-based AEADs should be used instead.
  \item \textbf{No formal audit.} The codebase has not undergone an
        independent third-party security audit. It relies on well-reviewed
        libraries (libsodium, OpenSSL, libargon2) and the CRYSTALS-Kyber
        reference implementation, but the integration code itself is
        unaudited.
\end{itemize}

\subsection*{Recent security fixes (v1.0.4)}
\begin{itemize}[leftmargin=1.4em]
  \item Hardened the libsodium-backed secure entry buffer to zero the unused
        tail of its guarded allocation on every edit, so a typed-then-cleared
        password no longer persists in locked memory until the widget is freed.
  \item Fixed a potential integer overflow in the secure buffer's capacity
        growth that could wrap to zero on a pathological paste.
  \item Added \texttt{fsync} of the output file and its parent directory around
        the atomic rename, so a crash or power loss cannot publish a truncated
        or zero-length output file.
\end{itemize}

\subsection*{Recent security fixes (v1.0.3)}
\begin{itemize}[leftmargin=1.4em]
  \item Removed a use-after-free in the GUI progress timer that could fire
        against freed state if the window was closed immediately after a job
        finished.
  \item Prevented the output's temporary file from resolving to, and
        truncating, the input file.
  \item Tightened validation of Argon2id parameters from untrusted headers.
  \item Replaced silent truncation of over-long paths with an explicit error.
\end{itemize}

% ---------------------------------------------------------------------------
\section{Recommendations for users}

\begin{enumerate}[leftmargin=1.6em]
  \item Use a high-entropy passphrase; the KDF cannot compensate for a weak
        one.
  \item Choose \textbf{Medium} or \textbf{Strong} key strength for sensitive
        data.
  \item Enable \textbf{hybrid mode} for data that must stay confidential for
        years (post-quantum margin).
  \item Run Ciphers only on a trusted, up-to-date system; it cannot defend a
        compromised host.
  \item Keep an independent, secure backup of the password — there is no
        recovery path.
\end{enumerate}

\vfill
\begin{center}
\small\textcolor{ink}{\textit{Ciphers is released under the MIT License. This
report describes version 1.0.4 and is provided for informational purposes; it
is not a warranty of fitness for any particular purpose.}}
\end{center}

\end{document}
